\documentclass[11pt]{article}
\usepackage[a4paper,margin=26mm]{geometry}
\usepackage{amsmath,amssymb,mathtools}
\usepackage{unicode-math}
\setmainfont{Latin Modern Roman}
\setmathfont{Latin Modern Math}
\usepackage{xurl,hyperref}
\hypersetup{hidelinks,pdftitle={Problem statement}}
\newcommand{\sref}[2]{\hyperlink{source:#1:#2}{[S#2]}}
\newcommand{\cataloguescope}[1]{#1}
\begin{document}
\section{Rigidity and moduli — Question 3.1}\label{statement}

\subsection{Definitions and mathematical statement}

\paragraph{Shared notation.}$\mathbb N=\{1,2,\ldots\}$, and $\mathbb Z,\mathbb Q,\mathbb R,\mathbb C$ denote the integers, rationals, reals and complex numbers. Any other conventions are specified locally.


A codimension-one foliation $\mathcal F$ decomposes a $3$-manifold $M$ into surface leaves, locally as the horizontal slices of a product chart. In the terminology of the source, it is taut if a closed curve transverse to $\mathcal F$ meets every leaf. A minimal foliation means that every leaf is dense; a minimal surface means a surface with vanishing mean curvature for the specified Riemannian metric. These two uses of minimal are distinct. A manifold is atoroidal if it has no incompressible embedded torus other than a boundary-parallel one. Write $\widetilde M$ for its universal cover, and put a tilde on a foliation or lamination for its lift. For an integral $2$-homology class $u$, the Thurston norm $x(u)$ is the minimum of $\sum_j\max\{0,-\chi(S_j)\}$ over embedded oriented representatives $S=\bigcup_j S_j$; extend it by homogeneity and continuity to real homology. Its unit ball has a polyhedral structure. A fibered face is a face whose cone contains classes represented by fibers of fibrations over $S^1$. A pseudo-Anosov flow has stable and unstable hyperbolic directions off finitely many closed prong-singular orbits. Monotone equivalence allows collapsing foliated product intervals. Almost transversality to a pseudo-Anosov flow permits the usual modifications at singular orbits; geometric limits refer to limits of the foliated geometric structures.
\paragraph{Statement recorded in the supplied source.}
For atoroidal $M$, construct a polyhedron $\mathscr P_{\mathcal F}$ refining the Thurston norm ball so as to parameterize taut foliations with all of the following properties. Each open face $c$ determines a pseudo-Anosov flow $X_c$, or an appropriate equivalence class of such flows; associated foliations can be isotoped transverse or almost transverse to $X_c$, or have a sublamination monotonically equivalent into its stable or unstable singular foliation. A geometric limit of foliations associated to $c$ is associated to a face contained in $\overline c$. Finally there is a natural polyhedral map from $\mathscr P_{\mathcal F}$ to a subdivision of the Thurston norm unit ball. The construction includes a choice of the appropriate parameterization and flow equivalence relation, rather than presupposing a particular one. \sref{statement}{2}
\subsection{Short English statement}
Seek a polyhedral parameterization of all taut foliations, with flows assigned to faces, compatible limits on their boundaries, and a polyhedral map to the Thurston norm ball.
\subsection{Sources}
\begin{description}
\item[\hypertarget{source:statement:1}{S1}] ULAM.AI and the UnsolvedMath Contributors, \emph{UnsolvedMath: A Curated Collection of Open Mathematics Problems}, catalogue entry. CC BY 4.0. \url{https://huggingface.co/datasets/ulamai/UnsolvedMath}.
\item[\hypertarget{source:statement:2}{S2}] Danny Calegari, \emph{Problems in foliations and laminations of 3-manifolds}, version 0.78 (2002), Question 3.1; Definitions 1.1--1.10 and local definitions and remarks. \url{https://arxiv.org/abs/math/0209081v1}.
\end{description}
\label{end:statement}
\end{document}
